\section{Surprisal concentration for independent copies}
\label{sec:iid_surprisal}

Let $\rho$ be a density matrix on a finite-dimensional Hilbert space, with
nonnegative eigenvalues $(p_i)_i$ summing to one. Set
\begin{align}
 S&=-\sum_i p_i\log p_i,
 &V&=\sum_i p_i(-\log p_i-S)^2.
 \label{eq:iid_surprisal_moments}
\end{align}
Terms with $p_i=0$ are zero. For $k\ge0$, write
$\rho_k=\rho^{\otimes k}$ and $L_k=-\log\rho_k$, with the logarithm equal
to zero on the kernel. These results justify the independent-copy
concentration step in \cite[Section~7]{OpenAI2026AreaLaw}.

\begin{theorem}[The product spectral distribution]
\label{thm:iid_surprisal_spectral_distribution}
\lean{Matrix.PosSemidef.re_trace_finKronecker_mul_cfc_surprisal}
\leanok
For every real function $f$ and $k\ge0$,
\begin{align}
 \operatorname{Tr}(\rho_k f(L_k))
  &=\sum_{x=(x_0,\ldots,x_{k-1})}
       \left(\prod_jp_{x_j}\right)
       f\left(-\log\prod_jp_{x_j}\right).
 \label{eq:iid_surprisal_spectral_distribution}
\end{align}
This identity also holds for any positive semidefinite $\rho$ whose trace
need not be one. A real function on a finite spectrum defines its spectral
functional calculus without a global continuity assumption.
\end{theorem}
\begin{proof}
\leanok
Tensoring a unitary eigenbasis of $\rho$ gives a unitary eigenbasis of
$\rho_k$, with eigenvalues $\prod_jp_{x_j}$. Apply the spectral functional
calculus in that basis and use invariance of the trace under unitary
conjugation.
\end{proof}

\begin{theorem}[Finite-product Chebyshev bound]
\label{thm:iid_surprisal_chebyshev}
\lean{Entropy.surprisalTail_pi_le_variance}
\leanok
Let $(p_i)_i$ be any finite probability distribution, allowing zero weights,
and let $S,V$ be given by~(\ref{eq:iid_surprisal_moments}). For every $k\ge0$
and $w>0$,
\begin{align}
 \sum_{x:\,| -\log\prod_jp_{x_j}-kS |>w}\prod_jp_{x_j}
   &\le \frac{kV}{w^2}.
 \label{eq:iid_surprisal_chebyshev}
\end{align}
\end{theorem}
\begin{proof}
\leanok
Choose independent samples $X_0,\ldots,X_{k-1}$ with law $(p_i)_i$.
The sum $Y=\sum_j-\log p_{X_j}$ has expectation $kS$ and variance $kV$.
Chebyshev's inequality bounds the probability that $|Y-kS|\ge w$ by
$kV/w^2$. On every tuple of positive probability,
$Y=-\log\prod_jp_{X_j}$. The remaining tuples have probability zero and
contribute nothing to the displayed sum.
\end{proof}

\begin{theorem}[The entropy window of the actual density matrix]
\label{thm:iid_surprisal_density_tail}
\lean{Matrix.PosSemidef.re_trace_surprisalTail_finKronecker_le,Matrix.PosSemidef.re_trace_surprisalTail_finKronecker_three_quarters_le}
\leanok
\uses{thm:iid_surprisal_spectral_distribution,thm:iid_surprisal_chebyshev}
For every $k\ge0$ and $w>0$,
\begin{align}
 \operatorname{Tr}\!\left(\rho_k\mathbf1_{\{|L_k-kS|>w\}}\right)
   &\le \frac{kV}{w^2}.
 \label{eq:iid_surprisal_density_tail}
\end{align}
In particular, for $k>0$ and $w=k^{3/4}$,
\begin{align}
 \operatorname{Tr}\!\left(\rho_k\mathbf1_{\{|L_k-kS|>k^{3/4}\}}\right)
   &\le \frac{V}{\sqrt{k}}.
 \label{eq:iid_surprisal_rate}
\end{align}
Thus the mass outside this entropy window tends to zero for a fixed density
matrix as the number of copies tends to infinity.
\end{theorem}
\begin{proof}
\leanok
Apply~(\ref{eq:iid_surprisal_spectral_distribution}) to the indicator of
the complement of the window and then use~(\ref{eq:iid_surprisal_chebyshev}).
For $k>0$, substitution of $w=k^{3/4}$ gives
$kV/w^2=V/\sqrt{k}$.
\end{proof}

These concentration statements do not yet select an irreducible
representation label. That step also requires comparison with the label
observable and control of the number of labels.
